% Chapter 2: Potential Outcomes
% A First Course in Causal Inference - Peng Ding
% Rewritten in Stein motivation-first style with textbook formulas

\chapter{潜在结果框架}

\section{从相关到因果：为什么关联不够？}

在第一章中，我们见证了一个令人不安的现象：Yule-Simpson 悖论向我们展示了关联可以与因果相悖。两组数据分别显示的正向效应，在合并后可能变为负向——这警告我们，从关联推断因果是危险的。

但我们不禁要问：\textbf{关联毕竟是我们能观察到的，而因果是我们真正想知道的。那么，因果效应能否被清晰地定义，甚至被估计？}

这个问题引导我们走向一个优雅而强大的框架——\textbf{潜在结果框架（Potential Outcomes Framework）}，也被称为 \textbf{Neyman-Rubin 因果模型}。

这个框架的起源可以追溯到 Neyman 在 1923 年的开创性工作，以及 Rubin 在 1970 年代对其的进一步发展。它的核心思想其实出奇地简单：要谈论因果，我们必须想象（counterfactual）如果事情有所不同会发生什么。

\section{一个思想实验}

让我们从一个最简单的场景开始：假设我们想研究某种新药对患者康复的因果效应。

考虑第 $i$ 个患者。我们记：
\begin{itemize}
  \item $T_i = 1$：患者接受了治疗
  \item $T_i = 0$：患者未接受治疗（对照组）
\end{itemize}

现在，关键的问题来了：这个患者的结果会是什么？

在现实中，我们只能观察到一种情况——要么患者接受了治疗，要么没有。但为了谈论因果，我们必须同时考虑两种可能性：

\begin{itemize}
  \item $Y_i(1)$：\textbf{如果}患者接受治疗，会观察到的结果
  \item $Y_i(0)$：\textbf{如果}患者未接受治疗，会观察到的结果
\end{itemize}

这里，$Y_i(1)$ 和 $Y_i(0)$ 就是\textbf{潜在结果}——它们描述了在不同治疗状态下可能发生的结果，但我们永远无法同时观察到两者。

\begin{Remark}[注 2.1 — 什么是反事实]
\textbf{反事实}（counterfactual）指的是未被观测到的潜在结果。具体而言：
\begin{itemize}
  \item 对于治疗组个体 $i$（$T_i=1$），我们观察到 $Y_i(1)$，而 $Y_i(0)$ 是反事实
  \item 对于对照组个体 $i$（$T_i=0$），我们观察到 $Y_i(0)$，而 $Y_i(1)$ 是反事实
\end{itemize}
正因为潜在结果框架讨论的是``如果接受相反治疗会怎样''，所以它也被称为\textbf{反事实框架}（counterfactual framework）。
\end{Remark}

\userannotation{为什么叫``潜在''结果？}{因为它们描述的是``如果...会怎样''的情形，这种结果是``潜在的''而非``现实的''。这正是 counterfactual（反事实）思维的本质。}

\section{科学表（Science Table）}

\cite{rubin2005} 将潜在结果组成的 $n \times 2$ 矩阵称为\textbf{科学表（Science Table）}：

\begin{center}
\begin{tabular}{ccc}
\hline
$i$ & $Y_i(1)$ & $Y_i(0)$ \\
\hline
$1$ & $Y_1(1)$ & $Y_1(0)$ \\
$2$ & $Y_2(1)$ & $Y_2(0)$ \\
$\vdots$ & $\vdots$ & $\vdots$ \\
$n$ & $Y_n(1)$ & $Y_n(0)$ \\
\hline
\end{tabular}
\end{center}

这个表包含了所有我们需要知道的关于因果效应的信息——但问题是，我们永远只能观察到每一行的其中一个值。

\section{个体因果效应}

有了潜在结果，我们终于可以清晰地定义因果效应了。

对于个体 $i$，定义\textbf{个体因果效应（Individual Treatment Effect, ITE）}为：
\[
\tau_i = Y_i(1) - Y_i(0)
\]

这是治疗相对于对照的``净效果''。

但是，这里出现了一个根本性的困难。1923 年，Jerzy Neyman 在一篇关于田间试验的论文中首次明确指出了这个问题：

\begin{center}
\textit{对于每个个体，我们只能观察到一种潜在结果——实际发生的那种。另一种潜在结果是潜在的、永不可见的。}
\end{center}

这被称为\textbf{因果推断的根本问题（Fundamental Problem of Causal Inference）}。

\userannotation{为什么这个困难如此根本？}{因为科学上感兴趣的大多数因果问题，都涉及将现实与反事实进行比较。我们可以问``这种药对这个患者有效吗？''——但这个问题要求我们同时知道患者吃药和不吃药的结果，而这是不可能的。}

\section{平均因果效应}

既然个体因果效应不可测量，我们退而求其次，关注\textbf{平均因果效应（Average Causal Effect, ACE）}：
\[
\tau = \mathbb{E}[Y_i(1) - Y_i(0)] = \bar{Y}(1) - \bar{Y}(0)
\tag{2.1}\label{eq:ace}
\]

用更精确的有限总体记号：\footnote{推导见附录 \cref{sec:appendix-ace-finite-sample}}
\[
\tau = \frac{1}{n} \sum_{i=1}^n \{ Y_i(1) - Y_i(0) \} = \frac{1}{n} \sum_{i=1}^n Y_i(1) - \frac{1}{n} \sum_{i=1}^n Y_i(0)
\tag{2.2}\label{eq:ace-finite}
\]

\textbf{一个关键恒等式}：平均因果效应等于个体因果效应的均值：
\[
\tau = \bar{Y}(1) - \bar{Y}(0) = \frac{1}{n} \sum_{i=1}^n \tau_i
\tag{2.3}\label{eq:ace-ite-mean}
\]

\section{亚组与亚组因果效应}

\textbf{亚组（subgroup）}是由某个协变量（通常是二元变量 $X_i$）分割出来的子群体。例如：
\begin{itemize}
  \item 按性别分割：男性亚组（$X_i=1$）和女性亚组（$X_i=0$）
  \item 按年龄分割：老年人/年轻人
  \item 按收入分割：高收入/低收入
\end{itemize}

\textbf{亚组因果效应（Subgroup Causal Effect）}衡量的是在该亚组内的平均因果效应：\footnote{推导见附录 \cref{sec:appendix-subgroup-ace}}
\[
\tau_x = \frac{\sum_{i=1}^n \mathbb{I}(X_i = x) \{ Y_i(1) - Y_i(0) \}}{\sum_{i=1}^n \mathbb{I}(X_i = x)}, \quad (x = 0, 1)
\tag{2.4}\label{eq:subgroup-ace}
\]

\section{亚组因果效应与 Yule-Simpson 悖论的不存在}

一个重要的恒等式揭示了亚组因果效应与总体因果效应的关系：\footnote{推导见附录 \cref{sec:appendix-subgroup-decomposition}}
\[
\tau = \pi_1 \tau_1 + \pi_0 \tau_0
\tag{2.5}\label{eq:ace-subgroup-decomposition}
\]
其中 $\pi_x = \frac{1}{n} \sum_{i=1}^n \mathbb{I}(X_i = x)$ 是 $X_i = x$ 的单元比例。

\textbf{关键推论}：如果 $\tau_1 > 0$ 且 $\tau_0 > 0$，则必有 $\tau > 0$。因此，\textbf{Yule-Simpson 悖论在因果效应层面不可能发生！}

这澄清了第一章的困惑：在关联层面出现的悖论，在因果层面并不存在——因为因果效应是可分解的，而关联并非如此。

\section{治疗分配机制与基本恒等式}

设 $T_i$ 是单元 $i$ 的二元治疗指示变量，观测结果 $Y_i$ 是潜在结果和治疗分配的函数：

\begin{align}
Y_i &= \begin{cases}
Y_i(1), & \text{if } T_i = 1 \\
Y_i(0), & \text{if } T_i = 0
\end{cases} \label{eq:observed-y} \\
&= T_i Y_i(1) + (1 - T_i) Y_i(0) \label{eq:fundamental-bridge} \\
&= Y_i(0) + T_i \{ Y_i(1) - Y_i(0) \} \label{eq:treatment-effect} \\
&= Y_i(0) + T_i \tau_i \label{eq:treatment-effect-2}
\end{align}

\cite{pearl2010} 将 \cref{eq:fundamental-bridge} 视为潜在结果与观测结果之间的\textbf{基本桥梁}。这些方程强调了个体因果效应 $\tau_i = Y_i(1) - Y_i(0)$ 可能在不同单元间存在异质性。

\section{SUTVA：使框架有效的必要假设}

使用潜在结果框架需要一些技术性假设。其中最重要的是\textbf{稳定单位治疗值假设（Stable Unit Treatment Value Assumption, SUTVA）}。

\begin{Assumption}[无干扰]\label{assump:sutva1}
单元 $i$ 的潜在结果不依赖于其他单元的治疗分配。这有时被称为\textbf{无干扰}假设。
\end{Assumption}

\begin{Assumption}[一致性]\label{assump:sutva2}
治疗水平没有其他版本。等价地，我们要求治疗水平是良好定义的，至少对关注的结局是如此。这有时被称为\textbf{一致性}假设。
\end{Assumption}

\begin{Assumption}[SUTVA]\label{assump:sutva-complete}
假设 \ref{assump:sutva1} 和 \ref{assump:sutva2} 同时成立。
\end{Assumption}

\textbf{为什么 SUTVA 如此重要？} 如果单元之间存在干扰，或者治疗定义不明确，那么潜在结果 $Y_i(1)$ 和 $Y_i(0)$ 就不是well-defined的——它们的值将取决于其他人的行为或治疗的具体形式，整个框架就会崩塌。

\textbf{违反 SUTVA 的例子}：

\textbf{无干扰假设（no-interference assumption）违反的情形}：

\begin{enumerate}
  \item \textbf{传染病与疫苗接种（infectious diseases and vaccination）}：当疾病具有传染性时，一个人是否接种疫苗不仅影响他自己的患病风险，还影响他周围人被感染的概率。例如，如果我的同事接种了流感疫苗，我患流感的概率也会降低——即使我自己没有接种。在这种情况下，单元 $i$ 的潜在结果 $Y_i(1)$ 和 $Y_i(0)$ 都依赖于其他人的治疗分配。

  \item \textbf{网络效应实验（network experiments）}：在社交网络实验中，如果一个人收到 Facebook 广告（treatment），他的朋友看到这条广告并受到影响后，可能也会购买该产品——即使朋友本人没有直接收到广告。这就是所谓的\textbf{溢出效应（spillover effect）}。

  \item \textbf{随机区组实验中的相邻地块（neighboring plots in randomized block experiments）}：农业实验中，相邻地块的作物可能相互影响——如果相邻地块使用了肥料，可能会通过土壤渗透影响本地块的产量。
\end{enumerate}

\textbf{一致性假设（consistency assumption）违反的情形}：

\begin{enumerate}
  \item \textbf{复合治疗（composite treatments）}：当治疗本身有复杂组成部分时，相同的``治疗''标签可能对应不同的实际干预。例如：
  \begin{itemize}
    \item 研究吸烟（smoking）对肺癌的影响时，``吸烟''这个治疗有多种形式——香烟类型、吸烟频率、吸入深度等，这些都可能导致不同的潜在结果。
    \item 研究大学教育（college education）对收入的影响时，``上大学''这个治疗包括大学类型（985/211 vs 普通院校）、专业、文凭等级等，这些差异会造成潜在结果的不一致。
  \end{itemize}

  \item \textbf{剂量效应（dosage effects）}：在医学研究中，相同的药物名称可能对应不同剂量，而剂量直接影响疗效。例如``服用阿司匹林''作为治疗，但不同剂量的阿司匹林会产生完全不同的潜在结果。

  \item \textbf{实施变异（implementation variation）}：即使治疗名义上相同，不同实施者可能引入差异。例如，心理治疗研究中，不同治疗师的风格和技能可能导致相同的``心理治疗''产生不同效果。
\end{enumerate}

当 SUTVA 不成立时，潜在结果框架需要修正。\textbf{现代因果推断文献}（如 \cite{halloran1996}）正在积极研究如何处理存在干扰或治疗版本多样的情形。

\userannotation{SUTVA 的历史重要性}{Rubin (1980) 将 SUTVA 确立为潜在结果框架的基础假设，没有 SUTVA，因果推断的数学化是不可能的。}

\section{反事实思维的逻辑}

也许你会好奇：为什么不能简单地问``实际发生的结果是什么''？

答案是：\textbf{因果效应本质上是一个反事实（counterfactual）概念}。当我们问``这种药是否有效''时，我们实际上是在问：

\begin{center}
\textit{``如果这个患者吃了药，结果会怎样？但现实是，他可能没有吃药...''}
\end{center}

这种思维方式可能让人觉得抽象甚至哲学化，但它是我们谈论因果的唯一严谨方式。潜在结果框架的价值在于，它将这种哲学洞察转化为清晰的数学语言。

\section{实验单元定义的微妙性}

我讨论一个与实验单元定义相关的微妙之处。简单来说，\textbf{实验单元可能不同于物理单元}。

例如，如果我在服药前没有服用阿司匹林且头痛没有消失，但我现在服用阿司匹林后头痛消失了，你可能认为我们可以同时观察到我的两种潜在结果。但这是对实验单元定义的误解！在不同时刻，我是同一个物理人，但成为了两个不同的实验单元。

因此，我们有四个潜在结果，其中两个被观察到，两个缺失。这说明，即使观察到的数据显示治疗``有效''，我们也无法确定地说阿司匹林就是原因——因为在``治疗后''的状态下，可能存在其他因素（如时间效应）导致头痛消失。

\section{非线性因果参数}

平均因果效应是一个\textbf{线性}因果参数，因为潜在结果的平均差等于个体差别的平均。其他参数可能不是线性的。

例如，中位数治疗效应：
\[
\delta_1 = \text{median}\{ Y_i(1) \}_{i=1}^n - \text{median}\{ Y_i(0) \}_{i=1}^n
\tag{2.6}\label{eq:median-treatment-effect-1}
\]
通常不同于个体治疗效应的中位数：
\[
\delta_2 = \text{median}\{ Y_i(1) - Y_i(0) \}_{i=1}^n
\tag{2.7}\label{eq:median-treatment-effect-2}
\]

这两个量在某些情况下可能相差很大。\footnote{数值例子见附录 \cref{sec:appendix-nonlinear-causal}}

\section{本章小结}

本章建立因果推断的数学语言：

\begin{enumerate}
  \item \textbf{潜在结果}：为每个个体定义 $Y_i(1)$ 和 $Y_i(0)$
  \item \textbf{科学表}：$n \times 2$ 矩阵 $\{Y_i(1), Y_i(0)\}_{i=1}^n$
  \item \textbf{个体因果效应}：$\tau_i = Y_i(1) - Y_i(0)$
  \item \textbf{平均因果效应}：$\tau = \bar{Y}(1) - \bar{Y}(0) = n^{-1}\sum_{i=1}^n \tau_i$
  \item \textbf{根本问题}：我们只能观察到一种潜在结果
  \item \textbf{基本桥梁方程}：$Y_i = T_i Y_i(1) + (1-T_i) Y_i(0)$
  \item \textbf{SUTVA}：一致性 + 无干扰
  \item \textbf{Yule-Simpson 悖论}在因果框架下不可能发生
\end{enumerate}

下一章我们将看到，\textbf{随机实验}提供了一种绕过根本问题的优雅方式——通过随机分配治疗，我们可以建立组间可比性，从而识别因果效应。

\section{习题}\label{sec:chapter2-exercises}

\begin{Exercise}{\ref{ex:2.1} 完美的医生}\label{ex:2.1}
假设潜在结果由以下模型生成：
\[
Y_i(0) \sim N(0, 1), \quad \tau_i = -0.5 + Y_i(0), \quad Y_i(1) = Y_i(0) + \tau_i
\]
治疗分配为 $T_i = \mathbb{I}(\tau_i \geq 0)$，观测结果为 $Y_i = T_i Y_i(1) + (1-T_i) Y_i(0)$。

计算条件均值差 $\mathbb{E}(Y_i \mid T_i=1) - \mathbb{E}(Y_i \mid T_i=0)$。

\textbf{提示}：需要使用截断正态分布的均值公式。当 $X \sim N(\mu, \sigma^2)$ 时，
\[
\mathbb{E}(X \mid a < X < b) = \mu - \sigma \frac{\phi\left(\frac{b-\mu}{\sigma}\right) - \phi\left(\frac{a-\mu}{\sigma}\right)}{\Phi\left(\frac{b-\mu}{\sigma}\right) - \Phi\left(\frac{a-\mu}{\sigma}\right)}
\]
其中 $\phi$ 和 $\Phi$ 分别是标准正态分布的密度函数和分布函数。
\end{Exercise}

\begin{Exercise}{\ref{ex:2.2} 非线性因果参数}\label{ex:2.2}
证明平均因果效应是线性因果参数的例子，即 $\bar{Y}(1) - \bar{Y}(0) = n^{-1}\sum_{i=1}^n \{Y_i(1) - Y_i(0)\}$。

其他参数可能不是线性的。例如，定义中位数治疗效应为：
\[
\delta_1 = \text{median}\{Y_i(1)\} - \text{median}\{Y_i(0)\}
\]
这通常不同于个体治疗效应的中位数：
\[
\delta_2 = \text{median}\{Y_i(1) - Y_i(0)\}
\]

\begin{enumerate}
  \item 给出数值例子分别满足 $\delta_1 = \delta_2$、$\delta_1 > \delta_2$ 和 $\delta_1 < \delta_2$。
  \item 哪个参数更有意义？为什么？请用例子来证明你的结论。
\end{enumerate}
\end{Exercise}

\begin{Exercise}{\ref{ex:2.3} 平均效应与个体效应}\label{ex:2.3}
给出一个数值例子，其中 $\tau = n^{-1}\sum_{i=1}^n \{Y_i(1) - Y_i(0)\} > 0$，但有不到一半的单元满足 $Y_i(1) > Y_i(0)$。即平均因果效应为正，但治疗对不到一半的单元有益。
\end{Exercise}

\begin{Exercise}{\ref{ex:2.4} 推荐阅读}\label{ex:2.4}
\begingroup
\parskip=0pt
\textbf{推荐阅读}：\cite{holland1986} 是统计因果推断领域的经典综述文章，它推广了``Rubin 因果模型''这一名称。在 UC Berkeley，我们称之为``Neyman 模型''——原因显而易见。
\endgroup
\end{Exercise}

